% --- Archon physics formalization source begin ---
% archon:physics
% archon:covers PhyXMiniProblems/problem_phyx_mini_0027.lean
% archon:source-report reports/phyx_mini/problem_phyx_mini_0027.source.json

\chapter{Physics problem phyx\_mini\_0027}
\label{ch:PhyXMiniProblems_problem_phyx_mini_0027}

\paragraph{Problem source.}
\#\# Physical scenario

In figure, a thin converging lens of focal length 14.0 cm forms an image of the square \textbackslash{}( abcd \textbackslash{}), which is \textbackslash{}( h\_c = h\_b = 10.0 \textbackslash{}text\{ cm\} \textbackslash{}) high and lies between distances of \textbackslash{}( p\_d = 20.0 \textbackslash{}text\{ cm\} \textbackslash{}) and \textbackslash{}( p\_a = 30.0 \textbackslash{}text\{ cm\} \textbackslash{}) from the lens. Let \textbackslash{}( a' \textbackslash{}), \textbackslash{}( b' \textbackslash{}), \textbackslash{}( c' \textbackslash{}), and \textbackslash{}( d' \textbackslash{}) represent the respective corners of the image. Let \textbackslash{}( q\_a \textbackslash{}) represent the image distance for points \textbackslash{}( a' \textbackslash{}) and \textbackslash{}( b' \textbackslash{}), \textbackslash{}( q\_d \textbackslash{}) represent the image distance for points \textbackslash{}( c' \textbackslash{}) and \textbackslash{}( d' \textbackslash{}),\textbackslash{}( h\_b' \textbackslash{}) represent the distance from point \textbackslash{}( b' \textbackslash{}) to the axis, and \textbackslash{}( h\_c' \textbackslash{}) represent the height of \textbackslash{}( c' \textbackslash{}).

\#\# Question

Find \textbackslash{}( h\_c' \textbackslash{}).

\#\# Answer choices

- A: A: \textbackslash{}( -27.1 \textbackslash{}text\{ cm\} \textbackslash{})

- B: B: \textbackslash{}( -19.6 \textbackslash{}text\{ cm\} \textbackslash{})

- C: C: \textbackslash{}( -18.5 \textbackslash{}text\{ cm\} \textbackslash{})

- D: D: \textbackslash{}( -23.3 \textbackslash{}text\{ cm\} \textbackslash{})

\#\# Figure caption (auxiliary; use the image as primary evidence)

The image shows a diagram of a convex lens setup. Here's a breakdown:

- **Objects and Labels**:
  - A vertical lens in the center with a double-headed arrow on top indicating its orientation.
  - Two focal points labeled "F" on the principal axis, one on each side of the lens.
  - A small rectangular object labeled with letters: "a" at the bottom-left, "b" at the top-left, "d" at the bottom-right, and "c" at the top-right, situated to the left of the lens.

- **Text and Measurements**:
  - The distance from the object to the lens is marked as "pₐ" with a double-headed arrow pointing towards the lens from the object.

- **Lines and Relationships**:
  - The principal axis is a horizontal line passing through the focal points and the center of the lens.
  - The object is upright and perpendicular to this axis.

The setup represents a typical ray diagram used in optics to illustrate image formation by a convex lens.

\#\# Dataset metadata

Geometrical Optics; Spatial Relation Reasoning, Multi-Formula Reasoning

\paragraph{Recorded answer/context.}
D: D: \textbackslash{}( -23.3 \textbackslash{}text\{ cm\} \textbackslash{})

\paragraph{Figure/image path.}
/root/proposal\_for\_physic/science-mango/phyx\_mini\_run/phyx\_data/test\_image/27.png

\paragraph{Formalization target.}
create a compiling Lean file with sorry bodies at `PhyXMiniProblems/problem\_phyx\_mini\_0027.lean`. The Lean declarations must preserve the physical quantities, dimensions or dimensional roles, figure labels, governing-law hypotheses, and final relation expressed by this problem.
Use Mathlib/Physlib names found through LeanExplore where available. If a physics API is missing, introduce faithful local abstractions rather than scalar placeholder aliases.

\begin{theorem}[Physics formalization target]
\label{thm:physics:phyx_mini_0027:target}
\lean{PhyXMiniProblems.ProblemPhyXMini0027.problem_phyx_mini_0027}
\uses{def:physics:phyx-mini-0027:phyxminiproblems-problemphyxmini0027-lengthincentimeters, def:physics:phyx-mini-0027:phyxminiproblems-problemphyxmini0027-thinlenssquaresetup, def:physics:phyx-mini-0027:phyxminiproblems-problemphyxmini0027-hasphysicalrealimageconfiguration, def:physics:phyx-mini-0027:phyxminiproblems-problemphyxmini0027-matchesfigurereadouts, def:physics:phyx-mini-0027:phyxminiproblems-problemphyxmini0027-respectslabeledimageplanes, def:physics:phyx-mini-0027:phyxminiproblems-problemphyxmini0027-satisfiesthinlensequation, def:physics:phyx-mini-0027:phyxminiproblems-problemphyxmini0027-satisfiestransversemagnification, def:physics:phyx-mini-0027:phyxminiproblems-problemphyxmini0027-matchesanswertonearesttenth}
The assigned autoformalize agent should translate this physics problem into Lean declarations in the covered file, with theorem and lemma proof bodies written as `by sorry`.
\end{theorem}
\begin{proof}
This is an autoformalization task, not a proof task. Produce statements that can later be proved by the physics prover without weakening the physical model.
\end{proof}

% --- Archon declaration topology begin ---
\section{Declaration topology}
\begin{definition}[Length Quantity]
  \label{def:physics:phyx-mini-0027:phyxminiproblems-problemphyxmini0027-lengthquantity}
  \lean{PhyXMiniProblems.ProblemPhyXMini0027.LengthQuantity}
  A signed physical length, independent of the unit used to read it
\end{definition}
\begin{proof}
  This is the declared model definition of Length Quantity.
\end{proof}
\begin{definition}[length In Centimeters]
  \label{def:physics:phyx-mini-0027:phyxminiproblems-problemphyxmini0027-lengthincentimeters}
  \lean{PhyXMiniProblems.ProblemPhyXMini0027.lengthInCentimeters}
  \uses{def:physics:phyx-mini-0027:phyxminiproblems-problemphyxmini0027-lengthquantity}
  The scalar centimeter readout of a signed physical length
\end{definition}
\begin{proof}
  This is the declared model definition of length In Centimeters.
\end{proof}
\begin{definition}[Object Corner]
  \label{def:physics:phyx-mini-0027:phyxminiproblems-problemphyxmini0027-objectcorner}
  \lean{PhyXMiniProblems.ProblemPhyXMini0027.ObjectCorner}
  The four object corners labeled in the figure
\end{definition}
\begin{proof}
  This is the declared model definition of Object Corner.
\end{proof}
\begin{definition}[Image Corner]
  \label{def:physics:phyx-mini-0027:phyxminiproblems-problemphyxmini0027-imagecorner}
  \lean{PhyXMiniProblems.ProblemPhyXMini0027.ImageCorner}
  The four corresponding image corners labeled with primes in the problem
\end{definition}
\begin{proof}
  This is the declared model definition of Image Corner.
\end{proof}
\begin{definition}[corresponding Image]
  \label{def:physics:phyx-mini-0027:phyxminiproblems-problemphyxmini0027-correspondingimage}
  \lean{PhyXMiniProblems.ProblemPhyXMini0027.correspondingImage}
  \uses{def:physics:phyx-mini-0027:phyxminiproblems-problemphyxmini0027-objectcorner, def:physics:phyx-mini-0027:phyxminiproblems-problemphyxmini0027-imagecorner}
  The correspondence between each object corner and its geometrical image
\end{definition}
\begin{proof}
  This is the declared model definition of corresponding Image.
\end{proof}
\begin{definition}[Thin Lens Square Setup]
  \label{def:physics:phyx-mini-0027:phyxminiproblems-problemphyxmini0027-thinlenssquaresetup}
  \lean{PhyXMiniProblems.ProblemPhyXMini0027.ThinLensSquareSetup}
  \uses{def:physics:phyx-mini-0027:phyxminiproblems-problemphyxmini0027-lengthquantity, def:physics:phyx-mini-0027:phyxminiproblems-problemphyxmini0027-objectcorner, def:physics:phyx-mini-0027:phyxminiproblems-problemphyxmini0027-imagecorner}
  The dimensionful quantities in the thin-lens ray diagram. objectDistance and imageDistance are positive distances measured away from the lens on their respective sides. objectHeight and imageHeight are signed transverse coordinates measured from the principal axis
\end{definition}
\begin{proof}
  This is the declared model definition of Thin Lens Square Setup.
\end{proof}
\begin{definition}[Has Physical Real Image Configuration]
  \label{def:physics:phyx-mini-0027:phyxminiproblems-problemphyxmini0027-hasphysicalrealimageconfiguration}
  \lean{PhyXMiniProblems.ProblemPhyXMini0027.HasPhysicalRealImageConfiguration}
  \uses{def:physics:phyx-mini-0027:phyxminiproblems-problemphyxmini0027-lengthincentimeters, def:physics:phyx-mini-0027:phyxminiproblems-problemphyxmini0027-correspondingimage, def:physics:phyx-mini-0027:phyxminiproblems-problemphyxmini0027-thinlenssquaresetup}
  The lens is converging and each pictured object point lies beyond its focal plane, producing a real image on the opposite side of the lens
\end{definition}
\begin{proof}
  This is the declared model definition of Has Physical Real Image Configuration.
\end{proof}
\begin{definition}[Matches Figure Readouts]
  \label{def:physics:phyx-mini-0027:phyxminiproblems-problemphyxmini0027-matchesfigurereadouts}
  \lean{PhyXMiniProblems.ProblemPhyXMini0027.MatchesFigureReadouts}
  \uses{def:physics:phyx-mini-0027:phyxminiproblems-problemphyxmini0027-lengthincentimeters, def:physics:phyx-mini-0027:phyxminiproblems-problemphyxmini0027-thinlenssquaresetup}
  Numerical and geometrical readouts supplied by the problem and figure: f = 14 cm, the far edge ab is at 30 cm, the near edge dc is at 20 cm, and the square extends from the axis to height 10 cm
\end{definition}
\begin{proof}
  This is the declared model definition of Matches Figure Readouts.
\end{proof}
\begin{definition}[Respects Labeled Image Planes]
  \label{def:physics:phyx-mini-0027:phyxminiproblems-problemphyxmini0027-respectslabeledimageplanes}
  \lean{PhyXMiniProblems.ProblemPhyXMini0027.RespectsLabeledImagePlanes}
  \uses{def:physics:phyx-mini-0027:phyxminiproblems-problemphyxmini0027-thinlenssquaresetup}
  The names q\_a and q\_d in the statement each denote one image plane: a', b' share q\_a, while c', d' share q\_d
\end{definition}
\begin{proof}
  This is the declared model definition of Respects Labeled Image Planes.
\end{proof}
\begin{definition}[Satisfies Thin Lens Equation]
  \label{def:physics:phyx-mini-0027:phyxminiproblems-problemphyxmini0027-satisfiesthinlensequation}
  \lean{PhyXMiniProblems.ProblemPhyXMini0027.SatisfiesThinLensEquation}
  \uses{def:physics:phyx-mini-0027:phyxminiproblems-problemphyxmini0027-lengthincentimeters, def:physics:phyx-mini-0027:phyxminiproblems-problemphyxmini0027-correspondingimage, def:physics:phyx-mini-0027:phyxminiproblems-problemphyxmini0027-thinlenssquaresetup}
  The governing thin-lens equation 1/f = 1/p + 1/q, applied to every corner. Writing it in one fixed length unit preserves the physical relation
\end{definition}
\begin{proof}
  This is the declared model definition of Satisfies Thin Lens Equation.
\end{proof}
\begin{definition}[Satisfies Transverse Magnification]
  \label{def:physics:phyx-mini-0027:phyxminiproblems-problemphyxmini0027-satisfiestransversemagnification}
  \lean{PhyXMiniProblems.ProblemPhyXMini0027.SatisfiesTransverseMagnification}
  \uses{def:physics:phyx-mini-0027:phyxminiproblems-problemphyxmini0027-lengthincentimeters, def:physics:phyx-mini-0027:phyxminiproblems-problemphyxmini0027-correspondingimage, def:physics:phyx-mini-0027:phyxminiproblems-problemphyxmini0027-thinlenssquaresetup}
  The signed transverse-magnification law h\_i / h\_o = -q / p, written in a division-free form that also applies to the two corners on the axis
\end{definition}
\begin{proof}
  This is the declared model definition of Satisfies Transverse Magnification.
\end{proof}
\begin{definition}[Answer Choice]
  \label{def:physics:phyx-mini-0027:phyxminiproblems-problemphyxmini0027-answerchoice}
  \lean{PhyXMiniProblems.ProblemPhyXMini0027.AnswerChoice}
  The four signed image-height choices printed in centimeters
\end{definition}
\begin{proof}
  This is the declared model definition of Answer Choice.
\end{proof}
\begin{definition}[answer Height In Centimeters]
  \label{def:physics:phyx-mini-0027:phyxminiproblems-problemphyxmini0027-answerheightincentimeters}
  \lean{PhyXMiniProblems.ProblemPhyXMini0027.answerHeightInCentimeters}
  \uses{def:physics:phyx-mini-0027:phyxminiproblems-problemphyxmini0027-answerchoice}
  The scalar centimeter readout displayed beside each answer choice
\end{definition}
\begin{proof}
  This is the declared model definition of answer Height In Centimeters.
\end{proof}
\begin{definition}[Matches Answer To Nearest Tenth]
  \label{def:physics:phyx-mini-0027:phyxminiproblems-problemphyxmini0027-matchesanswertonearesttenth}
  \lean{PhyXMiniProblems.ProblemPhyXMini0027.MatchesAnswerToNearestTenth}
  \uses{def:physics:phyx-mini-0027:phyxminiproblems-problemphyxmini0027-lengthquantity, def:physics:phyx-mini-0027:phyxminiproblems-problemphyxmini0027-lengthincentimeters, def:physics:phyx-mini-0027:phyxminiproblems-problemphyxmini0027-answerchoice, def:physics:phyx-mini-0027:phyxminiproblems-problemphyxmini0027-answerheightincentimeters}
  Agreement with an answer height displayed to the nearest tenth centimeter
\end{definition}
\begin{proof}
  This is the declared model definition of Matches Answer To Nearest Tenth.
\end{proof}
\begin{lemma}[named Image Distances eq]
  \label{lem:physics:phyx-mini-0027:phyxminiproblems-problemphyxmini0027-namedimagedistances-eq}
  \lean{PhyXMiniProblems.ProblemPhyXMini0027.namedImageDistances_eq}
  \uses{def:physics:phyx-mini-0027:phyxminiproblems-problemphyxmini0027-lengthincentimeters, def:physics:phyx-mini-0027:phyxminiproblems-problemphyxmini0027-thinlenssquaresetup, def:physics:phyx-mini-0027:phyxminiproblems-problemphyxmini0027-hasphysicalrealimageconfiguration, def:physics:phyx-mini-0027:phyxminiproblems-problemphyxmini0027-matchesfigurereadouts, def:physics:phyx-mini-0027:phyxminiproblems-problemphyxmini0027-respectslabeledimageplanes, def:physics:phyx-mini-0027:phyxminiproblems-problemphyxmini0027-satisfiesthinlensequation}
  The thin-lens equation determines the two named image distances q\_a = 105/4 cm and q\_d = 140/3 cm
\end{lemma}
\begin{proof}
  The conclusion follows from the governing relations and figure readouts named in the statement of named Image Distances eq.
\end{proof}
\begin{lemma}[image Height b Prime eq]
  \label{lem:physics:phyx-mini-0027:phyxminiproblems-problemphyxmini0027-imageheight-bprime-eq}
  \lean{PhyXMiniProblems.ProblemPhyXMini0027.imageHeight_bPrime_eq}
  \uses{def:physics:phyx-mini-0027:phyxminiproblems-problemphyxmini0027-lengthincentimeters, def:physics:phyx-mini-0027:phyxminiproblems-problemphyxmini0027-thinlenssquaresetup, def:physics:phyx-mini-0027:phyxminiproblems-problemphyxmini0027-hasphysicalrealimageconfiguration, def:physics:phyx-mini-0027:phyxminiproblems-problemphyxmini0027-matchesfigurereadouts, def:physics:phyx-mini-0027:phyxminiproblems-problemphyxmini0027-respectslabeledimageplanes, def:physics:phyx-mini-0027:phyxminiproblems-problemphyxmini0027-satisfiesthinlensequation, def:physics:phyx-mini-0027:phyxminiproblems-problemphyxmini0027-satisfiestransversemagnification}
  The far upper corner has signed image height h\_b' = -35/4 cm
\end{lemma}
\begin{proof}
  The conclusion follows from the governing relations and figure readouts named in the statement of image Height b Prime eq.
\end{proof}
\begin{lemma}[image Height c Prime eq]
  \label{lem:physics:phyx-mini-0027:phyxminiproblems-problemphyxmini0027-imageheight-cprime-eq}
  \lean{PhyXMiniProblems.ProblemPhyXMini0027.imageHeight_cPrime_eq}
  \uses{def:physics:phyx-mini-0027:phyxminiproblems-problemphyxmini0027-lengthincentimeters, def:physics:phyx-mini-0027:phyxminiproblems-problemphyxmini0027-thinlenssquaresetup, def:physics:phyx-mini-0027:phyxminiproblems-problemphyxmini0027-hasphysicalrealimageconfiguration, def:physics:phyx-mini-0027:phyxminiproblems-problemphyxmini0027-matchesfigurereadouts, def:physics:phyx-mini-0027:phyxminiproblems-problemphyxmini0027-respectslabeledimageplanes, def:physics:phyx-mini-0027:phyxminiproblems-problemphyxmini0027-satisfiesthinlensequation, def:physics:phyx-mini-0027:phyxminiproblems-problemphyxmini0027-satisfiestransversemagnification}
  For the near upper corner, the governing laws give the exact signed height h\_c' = -70/3 cm
\end{lemma}
\begin{proof}
  The conclusion follows from the governing relations and figure readouts named in the statement of image Height c Prime eq.
\end{proof}
% --- Archon declaration topology end ---
% --- Archon physics formalization source end ---
